\documentclass[]{beamer}

% Font encoding and input
\usepackage[T1]{fontenc}
\usepackage[utf8]{inputenc}
\usepackage{microtype}

% Core packages
\usepackage{graphicx}
\usepackage{url}
\usepackage{listings}
\usepackage{booktabs}
\usepackage{xcolor}
\usepackage{amsmath}
\usepackage{amsfonts}
\usepackage{amssymb}

\usepackage[spanish]{babel} %use your language

\setbeamertemplate{bibliography item}[text]

\usetheme{UTN-BHI}

% Enhanced code style
\lstdefinestyle{myCStyle}{
  language=C,
  basicstyle=\ttfamily\footnotesize,
  keywordstyle=\color{blue}\bfseries,
  commentstyle=\color{green!60!black},
  stringstyle=\color{red},
  numberstyle=\tiny\color{gray},
  showspaces=false,
  showstringspaces=false,
  breaklines=true,
  numbers=left,
  frame=single,
  backgroundcolor=\color{gray!10},
  captionpos=b
}

\urldef{\email}\path|jiparraguirre@frbb.utn.edu.ar|
\title[Juego Simple: Pong]{Juego Simple: Pong}
\subtitle{Integración de conceptos (semanas 9 a 13), de la consola a los gráficos}

\author[J. Iparraguirre] {Javier Iparraguirre}

\institute{
Universidad Tecnológica Nacional\\
11 de abril 461, Bahía Blanca, Argentina\\
\email\\
\url{http://www.frbb.utn.edu.ar/}}

\date{\today}

\begin{document}

\begin{frame}
  \titlepage
\end{frame}

\begin{frame}
  \frametitle{Contenido}
  \tableofcontents
\end{frame}

\section{Objetivo del proyecto}
\begin{frame}
  \frametitle{Objetivo del proyecto}
  \begin{itemize}
    \item Construir un juego de \textbf{Pong} mínimo que \textbf{integre} los conceptos vistos entre las semanas 9 y 13.
    \item Usar \textbf{únicamente variables simples}: sin estructuras y sin punteros en la lógica del juego.
    \item Mostrar la \textbf{misma lógica} en dos versiones que se pueden leer lado a lado:
    \begin{itemize}
      \item Versión de \textbf{consola} (dibuja con \texttt{printf}).
      \item Versión \textbf{gráfica} con \textbf{SDL3} (dibuja rectángulos de color).
    \end{itemize}
  \end{itemize}
\end{frame}

\section{El juego: Pong}
\begin{frame}
  \frametitle{El juego: Pong}
  \begin{itemize}
    \item Controlás la paleta izquierda; la computadora controla la paleta derecha.
    \item La pelota rebota en las paredes y en las paletas. Si pasa a un jugador, el otro suma un punto.
    \item Gana el primero en llegar a \textbf{5 puntos}.
  \end{itemize}
  \vspace{1em}
  \textbf{Controles}
  \begin{itemize}
    \item \texttt{W} $\rightarrow$ mover la paleta hacia arriba.
    \item \texttt{S} $\rightarrow$ mover la paleta hacia abajo.
    \item \texttt{ESC} $\rightarrow$ salir del juego.
  \end{itemize}
\end{frame}

\section{Conceptos integrados}
\begin{frame}
  \frametitle{Conceptos integrados (semanas 9 a 13)}
  \begin{table}
  \centering
  \begin{tabular}{@{}cll@{}}
  \toprule
  \textbf{Semana} & \textbf{Concepto} & \textbf{Dónde aparece en el juego} \\
  \midrule
  9  & Bucle \texttt{while}       & Bucle principal y lectura del teclado \\
  10 & Bucles \texttt{for}        & Dibujar el campo de juego \\
  11 & Funciones con prototipos   & Todo el juego se divide en funciones \\
  12 & Recursión                  & Cuenta regresiva del inicio \\
  13 & \texttt{switch}            & Control por teclado (W / S / ESC) \\
  \bottomrule
  \end{tabular}
  \caption{Cada semana aporta una pieza del programa.}
  \end{table}
\end{frame}

\section{Estructura del programa}
\begin{frame}[fragile]
  \frametitle{Estructura: variables globales y prototipos (semana 11)}
  \lstset{style=myCStyle}
  \begin{lstlisting}
/* Estado global del juego (sin estructuras). */
int ball_x, ball_y;      /* posicion de la pelota  */
int ball_dx, ball_dy;    /* direccion (-1 / +1)    */
int left_paddle_y;       /* paleta del jugador     */
int right_paddle_y;      /* paleta de la computadora */
int player_score, computer_score;
int game_running;

/* Prototipos de funciones (semana 11). */
void setup_game(void);
void countdown(int n);
void draw_field(void);
void handle_input(void);
void move_ball(void);
  \end{lstlisting}
\end{frame}

\section{Semana 9: el bucle principal}
\begin{frame}[fragile]
  \frametitle{Semana 9: el bucle principal \texttt{while}}
  \lstset{style=myCStyle}
  \begin{lstlisting}
countdown(3);      /* cuenta regresiva (recursion) */
setup_game();

/* Bucle principal del juego (semana 9). */
while (game_running)
{
    draw_field();      /* dibujar   (semana 10) */
    handle_input();    /* teclado   (semana 13) */
    move_ball();
    move_computer();
    check_paddles();
    Sleep(120);        /* SDL_Delay(120) en SDL */
}
  \end{lstlisting}
\end{frame}

\section{Semana 10: dibujar con \texttt{for}}
\begin{frame}[fragile]
  \frametitle{Semana 10: dibujar el campo con bucles \texttt{for}}
  \lstset{style=myCStyle}
  \begin{lstlisting}
/* Recorre el campo fila por fila. */
for (row = 1; row < HEIGHT - 1; row++)
{
    for (col = 0; col < WIDTH; col++)
    {
        if (col == 0 || col == WIDTH - 1)
            printf("#");        /* paredes    */
        else if (col == ball_x && row == ball_y)
            printf("O");        /* pelota     */
        else
            printf(" ");
    }
    printf("\n");
}
  \end{lstlisting}
\end{frame}

\section{Semana 12: recursión}
\begin{frame}[fragile]
  \frametitle{Semana 12: recursión (cuenta regresiva)}
  \lstset{style=myCStyle}
  \begin{lstlisting}
/* Cuenta regresiva recursiva antes de empezar. */
void countdown(int n)
{
    if (n == 0)          /* caso base */
    {
        printf("GO!\n");
        return;
    }
    printf("%d...\n", n);
    Sleep(700);
    countdown(n - 1);    /* llamada recursiva */
}
  \end{lstlisting}
\end{frame}

\section{Semana 13: \texttt{switch} y teclado}
\begin{frame}[fragile]
  \frametitle{Semana 13: \texttt{switch} para el control por teclado}
  \lstset{style=myCStyle}
  \begin{lstlisting}
while (_kbhit())
{
    key = _getch();
    switch (key)            /* integracion (semana 13) */
    {
    case 'w': case 'W':
        if (left_paddle_y > 1) left_paddle_y--;
        break;
    case 's': case 'S':
        if (left_paddle_y + PADDLE_HEIGHT < HEIGHT - 1)
            left_paddle_y++;
        break;
    case 27:                /* ESC */
        game_running = 0;
        break;
    }
}
  \end{lstlisting}
\end{frame}

\section{Sin estructuras ni punteros}
\begin{frame}
  \frametitle{Una decisión de diseño: variables simples}
  \begin{itemize}
    \item Todo el estado del juego se guarda en \textbf{variables globales simples} (\texttt{int}), no en estructuras.
    \item La lógica del juego \textbf{no usa punteros}: es la misma que ya sabemos escribir.
    \item Así el foco queda en los conceptos de las semanas 9 a 13, no en temas nuevos.
    \item En la versión SDL, los únicos punteros/estructuras inevitables son los que la \textbf{propia librería} necesita (ventana, renderer, \texttt{SDL\_FRect}).
  \end{itemize}
\end{frame}

\section{Dos versiones: consola y SDL3}
\begin{frame}
  \frametitle{La misma lógica, dos versiones}
  \small
  \begin{table}
  \centering
  \begin{tabular}{@{}p{3.1cm}p{3.0cm}p{3.4cm}@{}}
  \toprule
  \textbf{Concepto} & \textbf{Consola} & \textbf{SDL3} \\
  \midrule
  Bucle principal   & \texttt{while (game\_running)} & igual \\
  Dibujo (for)      & \texttt{printf} de caracteres  & \texttt{draw\_cell()} + rectángulos \\
  Recursión         & \texttt{countdown()}           & igual \\
  Teclado (switch)  & \texttt{\_getch()}             & \texttt{event.key.scancode} \\
  Entrada           & \texttt{\_kbhit()}            & \texttt{SDL\_PollEvent()} \\
  Espera            & \texttt{Sleep(120)}           & \texttt{SDL\_Delay(120)} \\
  \bottomrule
  \end{tabular}
  \caption{El esqueleto es idéntico; solo cambia el dibujo y la entrada.}
  \end{table}
\end{frame}

\begin{frame}[fragile]
  \frametitle{Versión SDL3: mismo \texttt{switch}, otra fuente de teclas}
  \lstset{style=myCStyle}
  \begin{lstlisting}
while (SDL_PollEvent(&event))
{
    if (event.type == SDL_EVENT_KEY_DOWN)
    {
        switch (event.key.scancode)
        {
        case SDL_SCANCODE_W: /* ... */ break;
        case SDL_SCANCODE_S: /* ... */ break;
        case SDL_SCANCODE_ESCAPE:
            game_running = 0;
            break;
        }
    }
}
  \end{lstlisting}
\end{frame}

\section{Cómo compilar y jugar}
\begin{frame}
  \frametitle{Cómo compilar y jugar}
  \textbf{Versión de consola} (\texttt{super-simple-game/})
  \begin{itemize}
    \item Compilar \texttt{super-simple-game.cpp} con un compilador de C/C++ en Windows.
    \item Usa \texttt{conio.h} y \texttt{windows.h} (consola de Windows).
  \end{itemize}
  \vspace{0.8em}
  \textbf{Versión gráfica} (\texttt{super-simple-game-SDL/})
  \begin{enumerate}
    \item Descargar SDL3 (\emph{SDL3-devel-3.x.x-VC.zip}) y crear la carpeta \texttt{SDL3/} junto al proyecto.
    \item Abrir \texttt{super-simple-game-SDL.sln} en Visual Studio 2022.
    \item Elegir la plataforma \textbf{x64} y presionar \textbf{F5}.
  \end{enumerate}
\end{frame}

\section{Conclusión}
\begin{frame}
  \frametitle{Progresión del aprendizaje}
  \begin{itemize}
    \item \textbf{Consola:} integramos while, for, funciones, recursión y switch en un programa completo.
    \item \textbf{SDL3:} el \textbf{mismo} programa aprende a dibujarse en una ventana gráfica.
    \item Idea clave: los conceptos básicos alcanzan para un juego real; el paso a gráficos \textbf{no} cambia la lógica.
  \end{itemize}
\end{frame}

\begin{frame}
  \begin{center}
    \begin{beamercolorbox}[sep=8pt,center]{title}
      \Large{¡Gracias!}
    \end{beamercolorbox}
    \vspace{1em}
    \huge{¿Preguntas?}\\
    \vspace{1em}
    \small{\email}
  \end{center}
\end{frame}

\section{Referencias}
\begin{frame}[allowframebreaks]
  \frametitle{Referencias}
  \bibliographystyle{alpha}
  \bibliography{references}
\end{frame}

\end{document}
